\documentclass[11pt,letterpaper]{article}

\usepackage[utf8]{inputenc}
\usepackage[T1]{fontenc}
\usepackage[spanish,es-nodecimaldot]{babel}
\usepackage{lmodern}
\usepackage{microtype}
\usepackage{geometry}
\usepackage{booktabs}
\usepackage{longtable}
\usepackage{graphicx}
\usepackage{xcolor}
\usepackage{hyperref}
\usepackage{parskip}
\usepackage{enumitem}
\usepackage{fancyhdr}
\usepackage{lastpage}

\geometry{margin=2.25cm}
\definecolor{AtoyacBlue}{HTML}{174A6A}
\definecolor{AtoyacGreen}{HTML}{2F6F5E}
\definecolor{AtoyacGray}{HTML}{555555}
\hypersetup{colorlinks=true, linkcolor=AtoyacBlue, urlcolor=AtoyacBlue}
\setlength{\parindent}{0pt}
\setlength{\parskip}{0.65em}
\setlength{\headheight}{14pt}
\setlist[itemize]{leftmargin=1.35em, itemsep=0.2em, topsep=0.2em}
\setlist[enumerate]{leftmargin=1.5em, itemsep=0.25em, topsep=0.25em}
\pagestyle{fancy}
\fancyhf{}
\lhead{\textcolor{AtoyacGray}{Auditoria DENUE textil}}
\rhead{\textcolor{AtoyacGray}{Clasificacion productiva}}
\cfoot{\textcolor{AtoyacGray}{\thepage\ de \pageref{LastPage}}}

\newcommand{\mapfig}[3]{
\begin{figure}[p]
\centering
\includegraphics[width=0.92\textwidth,height=0.72\textheight,keepaspectratio]{#1}
\caption{#2. La clasificacion representa una prioridad de investigacion basada en actividad declarada, senales textuales y ubicacion. No constituye evidencia de descarga ni contaminacion atribuible al establecimiento.}
\label{#3}
\end{figure}
}

\begin{document}

{\centering
{\Huge\bfseries\color{AtoyacBlue} Guia de auditoria y clasificacion DENUE textil\par}
{\Large\bfseries Actividad productiva probable y prioridad de investigacion\par}
\vspace{0.8em}
{\large Servicio social Atoyac\par}
{\normalsize \today\par}
\vspace{1em}
}
\hrule
\vspace{1em}

\section{Proposito}

Este documento guia la revision de la nueva clasificacion productiva DENUE. El objetivo no es afirmar contaminacion, descargas o incumplimiento. El objetivo es construir un universo reproducible de establecimientos prioritarios para investigacion documental y, cuando corresponda, validacion en campo.

La tabla central es:

\begin{verbatim}
outputs/tables/denue_clasificacion_productiva/
denue_universo_textil_clasificado.csv
\end{verbatim}

El campo central es \texttt{flag\_estudio\_prioritario}, acompanado por \texttt{motivo\_flag\_estudio\_prioritario}.

Para el alcance operativo de este proyecto, el campo mas importante es \texttt{flag\_universo\_alcance\_proyecto}. Ese campo marca los registros que pertenecen al universo real de analisis ambiental: procesos humedos o tratamiento, lavado/deslavado/lavanderia, o mezclilla/jeans. Los demas registros siguen siendo textiles y se conservan como contexto trazable, pero no forman el universo principal del alcance.

\section{Que registros se auditan}

No todos los registros se auditan con la misma intensidad. La auditoria prioriza:

\begin{itemize}
  \item establecimientos con senales de lavado, deslavado, lavanderia, tenido, tintoreria, acabado o tratamiento;
  \item registros con senales de mezclilla o jeans;
  \item maquila productiva o fabricacion textil con incertidumbre;
  \item registros cercanos a hidrografia, especialmente a 250 m o menos;
  \item casos ambiguos, contradictorios o con baja confianza.
\end{itemize}

El Excel de trabajo es:

\begin{verbatim}
outputs/tables/auditoria_enriquecida/
auditoria_denue_textil_priorizada.xlsx
\end{verbatim}

\section{Columnas que debe llenar la persona auditora}

\begin{longtable}{p{4cm}p{10cm}}
\toprule
Columna & Como llenarla \\
\midrule
\texttt{categoria\_auditada} & Alta, media, baja, excluir, pendiente o revisar, segun evidencia. \\
\texttt{etapa\_productiva\_auditada} & Confirmar o corregir la etapa sugerida: lavado, deslavado, tintoreria, acabado, confeccion, comercio, etc. \\
\texttt{notas\_auditoria} & Justificacion breve: evidencia observada, duda, fuente usada o razon para excluir. \\
\texttt{fuente\_auditoria} & DENUE, Google Maps, busqueda web, padron, directorio, evidencia fotografica, visita de campo u otra fuente. \\
\texttt{decision\_estudio} & Incluir prioritario, incluir como contexto productivo, excluir, mantener pendiente, validar documentalmente o validar en campo. \\
\texttt{fecha\_auditoria} & Fecha de revision en formato AAAA-MM-DD. \\
\texttt{auditor\_responsable} & Nombre, iniciales o equipo responsable. \\
\bottomrule
\end{longtable}

\section{Como interpretar los flags}

\begin{itemize}
  \item \texttt{flag\_estudio\_prioritario}: senales especialmente relevantes para investigar procesos de tratamiento de prendas, mezclilla o textiles con posible uso de agua.
  \item \texttt{flag\_universo\_alcance\_proyecto}: registros dentro del alcance real del proyecto: procesos humedos, lavado/deslavado/lavanderia o mezclilla/jeans.
  \item \texttt{flag\_proceso\_humedo\_relevante}: senales de lavado, tintoreria, tenido, acabado o tratamiento.
  \item \texttt{flag\_lavado\_deslavado}: subconjunto de alta atencion dentro de procesos humedos; incluye lavado, deslavado, prelavado o lavanderia.
  \item \texttt{flag\_mezclilla\_jeans}: senales de mezclilla, denim, jeans o pantalon de mezclilla.
  \item \texttt{flag\_maquila\_productiva}: senales de maquila o confeccion productiva.
  \item \texttt{flag\_cercania\_hidrografia\_250m}: distancia a hidrografia menor o igual a 250 m. Aumenta prioridad de revision, pero no prueba descarga.
  \item \texttt{flag\_posible\_falso\_positivo}: registro que entro al universo textil pero puede no ser relevante para el estudio prioritario.
  \item \texttt{flag\_posible\_falso\_negativo}: registro con senales mixtas que no conviene excluir automaticamente.
\end{itemize}

\section{Resultados de control}

\begin{longtable}{lr}
\toprule
Indicador & Registros \\
\midrule
Universo textil inicial & 1067 \\
Universo textil depurado & 1049 \\
Universo de alcance del proyecto & 327 \\
Universo ambientalmente prioritario & 266 \\
Pendientes de auditoria & 452 \\
Procesos humedos o tratamiento & 222 \\
Senales de lavado o deslavado & 171 \\
Senales de tenido o tintoreria & 20 \\
Senales de acabado o tratamiento & 5 \\
Senales de mezclilla o jeans & 110 \\
Maquila o confeccion productiva & 561 \\
Comercio o renta & 7 \\
Confeccion especializada de baja relevancia & 17 \\
Prioritarios a 250 m o menos de hidrografia & 83 \\
\bottomrule
\end{longtable}

\section{Fuentes recomendadas para auditar}

\begin{itemize}
  \item DENUE y campos originales del registro.
  \item Google Maps y Street View cuando exista cobertura.
  \item Busqueda web del nombre comercial.
  \item Padrones municipales, licencias, directorios o camaras empresariales.
  \item Evidencia fotografica o visita de campo.
  \item Informacion comunitaria o institucional disponible.
\end{itemize}

\section{Ejemplos de decision}

\begin{longtable}{p{4cm}p{4cm}p{6cm}}
\toprule
Caso & Decision sugerida & Criterio \\
\midrule
Lavanderia, tintoreria, deslavado o acabado & Incluir prioritario & Senal compatible con proceso humedo o tratamiento. \\
Jeans o mezclilla con maquila/fabricacion & Validar documentalmente o incluir prioritario si hay tratamiento & Puede formar parte de cadena productiva relevante. \\
Confecciones genericas & Contexto productivo o pendiente & No asumir proceso humedo sin evidencia adicional. \\
Boutique, tienda o comercio & Excluir del universo prioritario & Conservar trazabilidad, sin inflar el universo ambiental. \\
Novias, modista, sastreria o arreglos & Excluir del universo prioritario & Confeccion especializada de baja relevancia si no hay senales industriales. \\
Registro ambiguo cerca de hidrografia & Auditoria documental obligatoria & Cercania aumenta prioridad de revision, no prueba contaminacion. \\
\bottomrule
\end{longtable}

\section{Mapas clave}

\mapfig{../outputs/maps/denue_clasificacion_productiva/01_universo_textil_inicial.png}{Universo textil detectado inicialmente}{fig:universo-inicial}
\mapfig{../outputs/maps/denue_clasificacion_productiva/02_universo_textil_depurado.png}{Universo textil depurado}{fig:universo-depurado}
\mapfig{../outputs/maps/denue_clasificacion_productiva/03_establecimientos_prioritarios_estudio.png}{Establecimientos prioritarios para investigacion}{fig:prioritarios}
\mapfig{../outputs/maps/denue_clasificacion_productiva/04_procesos_humedos_relevantes.png}{Senales compatibles con procesos humedos o tratamiento}{fig:procesos-humedos}
\mapfig{../outputs/maps/denue_clasificacion_productiva/05_lavado_deslavado.png}{Senales de lavado y deslavado}{fig:lavado}
\mapfig{../outputs/maps/denue_clasificacion_productiva/07_tenido_tintoreria.png}{Senales de tenido y tintoreria}{fig:tenido}
\mapfig{../outputs/maps/denue_clasificacion_productiva/08_acabado_tratamiento.png}{Senales de acabado y tratamiento especial}{fig:acabado}
\mapfig{../outputs/maps/denue_clasificacion_productiva/09_mezclilla_jeans.png}{Senales de mezclilla y jeans}{fig:mezclilla}
\mapfig{../outputs/maps/denue_clasificacion_productiva/11_pendientes_auditoria.png}{Registros pendientes de auditoria}{fig:pendientes}
\mapfig{../outputs/maps/denue_clasificacion_productiva/14_prioritarios_250m_hidrografia.png}{Prioritarios a 250 m o menos de hidrografia}{fig:prioritarios-250}
\mapfig{../outputs/maps/denue_clasificacion_productiva/17_prioridad_campo.png}{Prioridad de campo o revision territorial}{fig:prioridad-campo}
\mapfig{../outputs/maps/denue_clasificacion_productiva/18_final_universo_recomendado.png}{Universo recomendado consolidado}{fig:final}
\mapfig{../outputs/maps/denue_clasificacion_productiva/19_universo_alcance_proyecto.png}{Universo de alcance del proyecto: procesos humedos, lavado/deslavado o mezclilla/jeans}{fig:alcance-proyecto}

\section{Mapas interactivos}

Los mapas interactivos Plotly se encuentran en:

\begin{verbatim}
outputs/maps_interactive/denue_clasificacion_productiva/index.html
\end{verbatim}

\section{Aplicacion de auditoria}

Despues de llenar el Excel, ejecutar:

\begin{verbatim}
python scripts/07_apply_denue_audit.py
\end{verbatim}

Ese script conserva compatibilidad con las plantillas antiguas y ademas lee el Excel priorizado para producir salidas auditadas de la clasificacion productiva.

\end{document}
